
\medskip
\noindent
\SGARefTarget{sga7.sga7i.eix.n9-6}{9.6}. We are now in a position to define the announced canonical pairing
(\SGARefLink{sga7.sga7i.eix.n9-1-2}{9.1.2}), without any restriction on $\ell$, in such a way that, for
$\ell\neq p$, this definition agrees with the one given in \SGARefLink{sga7.sga7i.eix.n9-2}{9.2}. For this
purpose, we apply the considerations of \SGARefLink{sga7.sga7i.eix.n9-5}{9.5} to the following situation:

\begin{equation}\tag{\SGARefTarget{sga7.sga7i.eix.n9-6-1}{9.6.1}}
\left\{
\begin{aligned}
P&=T_\ell(A^o)^t=\check{\underline M}{}'_\ell(1),
& Q&=\underline M_\ell,
& R&=T_\ell(A^o)^f/T_\ell(A^o)^t,\\[1.25ex]
F&=T_\ell(A^o)^f,
& E&=D(F')=D\bigl(T_\ell({A'}^o)^f\bigr).
\end{aligned}
\right.
\end{equation}
In the last formula, $D$ denotes Cartier duality in the sense of
pro-$\ell$ Barsotti--Tate groups; that is, $D(Y)$ is the
\underline{projective} system of the $D(Y(n))$, with the usual transition
morphisms. The structure of $F$ as an extension of $R$ by $P$ is clear.
Likewise,

\begin{equation*}
F'=T_\ell({A'}^o)^f
\end{equation*}
has the structure of an extension of $R'$ by
$P'=T_\ell({A'}^o)^t=\check{\underline M}_\ell(1)$. Applying Cartier
duality, we conclude that $E=D(F')$ is endowed with the structure of an
extension of $D(P')=\underline M_\ell$ by $D(R')$. Now $D(R')$ is
canonically isomorphic to $R$ by \SGARefLink{sga7.sga7i.eix.n5-5}{5.5} (proved in \SGARefLink{sga7.sga7i.eix.n7-4-6}{7.4.6}); this specifies the
structure of $E$ as an extension of $Q$ by $R$.

As the mixed extension $X_\eta$ of $E_\eta$ by $F_\eta$, we take

\begin{equation}\tag{\SGARefTarget{sga7.sga7i.eix.n9-6-2}{9.6.2}}
X_\eta=T_\ell(A_K)=X^0_\eta
 \supset X^1_\eta=T_\ell(A_K)^f
 \supset X^2_\eta=T_\ell(A_K)^t
 \supset0.
\end{equation}
The isomorphism of extensions of $R_\eta$ by $P_\eta$

\begin{equation*}
X^1_\eta\simeq F_\eta
\end{equation*}
is clear. The isomorphism of extensions of $Q_\eta$ by $R_\eta$

\begin{equation*}
X_\eta/X^2_\eta\simeq E_\eta=D(F'_\eta)
\end{equation*}
is clear from the orthogonality theorem \SGARefLink{sga7.sga7i.eix.n5-2}{5.2}.
\footnote{The printed formula has $F_\eta=D(F'_\eta)$. Since
$E=D(F')$ and the displayed quotient is the extension of $Q_\eta$ by
$R_\eta$, the required term is $E_\eta$.}

We therefore have the data required to define a pairing (\SGARefLink{sga7.sga7i.eix.n9-5-5}{9.5.5}), which
here, since $Q^*=\underline M'_\ell$, takes the form (\SGARefLink{sga7.sga7i.eix.n9-1-2}{9.1.2}).
\footnote{The printed source calls this a pairing ``(\SGARefLink{sga7.sga7i.eix.n9-5-6}{9.5.6})''; the pairing
is the homomorphism (\SGARefLink{sga7.sga7i.eix.n9-5-5}{9.5.5}), whereas \SGARefLink{sga7.sga7i.eix.n9-5-6}{9.5.6} is the subsequent proposition.}
Its compatibility with the definition of \SGARefLink{sga7.sga7i.eix.n9-2}{9.2} follows from \SGARefLink{sga7.sga7i.eix.n9-5-8}{9.5.8}. The
considerations of \SGARefLink{sga7.sga7i.eix.n5-8}{5.8} also show, subject to the validity of Tate's theorem
[33] (hence at least when $\operatorname{char}K=0$), that the pairing in
question can be described in terms of the single pro-$\ell$
Barsotti--Tate group $X_\eta=T_\ell(A_K)$ over $\eta$: its three-step
filtration, and the extensions $F$ and $E$ extending $X^1_\eta$ and
$X_\eta/X^2_\eta$, can be defined canonically in terms of $X_\eta$ alone.
\footnote{The printed sentence says that $X_\eta$ is ``over $S$''; its
notation and definition make it the generic-fibre group over $\eta$. It
also lacks its closing parenthesis.}

\medskip
\noindent
\underline{Remark \SGARefTarget{sga7.sga7i.eix.n9-7}{9.7}.} We no longer necessarily assume that $A_K$ has
semistable reduction over the henselian trait $S$. Let $K'$ be a finite
Galois extension of $K$ such that $A_{K'}$ has semistable reduction over
the normalization $S'$ of $S$ in $K'$ (\SGARefLink{sga7.sga7i.eix.n3-6}{3.6}). With the notation of (\SGARefLink{sga7.sga7i.eix.n8-1-8}{8.1.8}),
Galois descent from the monodromy pairing for $A_{K'}$ then defines a
pairing

\begin{equation}\tag{*}
v_\ell^{K'}:\underline M_K\otimes\underline M'_K
 \longrightarrow\mathbb Z_{\ell K}.
\end{equation}
A priori this pairing depends on the choice of $K'$, but it follows from
(10.4.3) below that the pairing

\begin{equation}\tag{\SGARefTarget{sga7.sga7i.eix.n9-7-1}{9.7.1}}
u_\ell:\underline M_K\otimes\underline M'_K
 \longrightarrow\mathbb Q_{\ell K}
\end{equation}
given by

\begin{equation*}
u_\ell=\frac{1}{e(S'/S)}\,v_\ell^{K'}
\end{equation*}
is independent of it, where

\begin{equation*}
e(S'/S)=\operatorname{length}(V'/\mathfrak mV').
\end{equation*}
We may again call (\SGARefLink{sga7.sga7i.eix.n9-7-1}{9.7.1}) the \underline{monodromy pairing} for the
abelian scheme $A_K$ over $K$ and the prime number $\ell$ under
consideration: when $A_K$ itself has semistable reduction, it plainly
coincides with the restriction to $\eta$ of the monodromy pairing (\SGARefLink{sga7.sga7i.eix.n9-1-2}{9.1.2}).
It follows from \SGARefLink{sga7.sga7i.eix.n10-4}{10.4} below that this pairing in fact always has values in
$\mathbb Q_K$ and is independent of $\ell$.

One may ask whether it has values in $\mathbb Z_K$, which is equivalent to
saying that (\SGARefLink{sga7.sga7i.eix.n9-7-1}{9.7.1}) has values in $\mathbb Z_{\ell K}$ for every prime
number $\ell$. This is so at least when $A_K$ is an elliptic curve. The
only case requiring verification is then the one in which
$\underline M_K\simeq\underline M'_K$ has rank $1$, so that
$\underline M_K\otimes\underline M'_K$ is canonically isomorphic to
$\mathbb Z_K$ and the monodromy pairing (\SGARefLink{sga7.sga7i.eix.n9-7-1}{9.7.1}) can simply be interpreted
as an ordinary rational number. It follows from the results of Néron [19]
that this number is $v(j)$, where $j$ is the absolute invariant of $A_K$
and $v$ is the normalized valuation of $K$.

